\section{Validacio es teszteles}
\label{ch:testing}

\subsection{Tesztelesi szintek}

A projektben a validacio tobb szinten ertelmezheto. Egyreszt igazolni kell, hogy a rendszer a definialt use case-ek alapjan mukodokepes. Masreszt technikai oldalrol is be kell mutatni, hogy a frontend, a backend es az API-kapcsolat stabilan egyutt tud mukodni. A fejezet javasolt szerkezete a kovetkezo:

\begin{itemize}
	\item frontend oldali komponens- vagy egysegtesztek,
	\item backend oldali API- vagy integracios ellenorzesek,
	\item manualis use case validacio,
	\item fejlesztoi es demonstracios kornyezetben vegzett vegponti kiprobalas.
\end{itemize}

\subsection{Meresi es ellenorzesi szempontok}

A kesobbi bovitett valtozatban itt erdemes rogzitni azokat a meresi szempontokat, amelyek alapjan a rendszer ertekelheto. Ilyen lehet peldaul a keresesi folyamat mukodese, a receptletrehozas sikeressege, a jogosultsaghoz kotott muveletek vedelme, illetve a komment- es ertekeleskezeles helyes viselkedese.

\subsection{Dokumentalhato bizonyitekok}

Ehhez a fejezethez jo alapot adhatnak a repositoryban mar meglevo fejlesztoi anyagok, peldaul a Bruno gyujtemenybe szervezett API-hivasok es a wireframe-ekhez kapcsolodo kepernyok. A vegleges dolgozatba erdemes majd tablazatos eredmenyeket, hibaeseteket es kepernyomentéseket is beemelni, hogy a validacio ne csak leiro jellegu legyen.